\section{Two-Phase Pretraining: cuándo deben aparecer los datos de alta calidad}

La segunda investigación sobre datos lleva la pregunta de la source selection a la schedule: si los tokens de alta calidad son escasos, ¿conviene aumentar su proporción al final del entrenamiento? El esquema de dos fases no consiste simplemente en «primero datos peores y luego mejores»; primero usa una broad mixture para construir capacidades generales y después usa una blend de alta calidad ya validada para modificar los gradientes de la etapa final, al mismo tiempo que evita la repetición excesiva.


\lecturefigure{slide-046.jpg}{Two-Phase Pretraining: datos organizados por fases para mejorar la LLM accuracy}{CS25 V5 official deck, p.~46}


\lecturefigure{slide-047.jpg}{Optimizar la LLM pretraining data selection: calidad, mezcla y schedule}{CS25 V5 official deck, p.~47}


\lecturefigure{slide-048.jpg}{Contribuciones principales: definición, evaluación, análisis fine-grained y escalabilidad}{CS25 V5 official deck, p.~48}

\teachervoice{00:21:03--00:22:30, Steven describe el trabajo como una validación sistemática del two-phase pretraining y enfatiza que primero se hacen prototipos de las blends con un token budget más pequeño. Sugerencia práctica: la receta de datos, al igual que los hiperparámetros del modelo, requiere cheap proxy experiments.}

\subsection{Formalización: phase y blend son dos controles independientes}

La sección anterior presentó la intuición de two-phase; esta sección la descompone en variables que pueden controlarse por separado. Solo con una formalización los investigadores pueden distinguir si la ganancia proviene de los datos de la segunda fase, del momento en que aparecen, de la duración del entrenamiento o de tokens adicionales. Esta distinción también determina el presupuesto experimental: fijar el total de tokens permite estudiar la schedule, fijar la phase duration permite estudiar la blend, y cambiar ambas a la vez solo produce una recipe a la que no puede atribuirse la causa.

Sea $q_1$ la distribución de datos de la primera fase con $T_1$ pasos, y $q_2$ la de la segunda fase con $T_2$ pasos; el objetivo total puede escribirse como
\[
\mathcal{L}=\sum_{t=1}^{T_1}\mathbb{E}_{x\sim q_1}[\ell_t(x)]
+\sum_{t=T_1+1}^{T_1+T_2}\mathbb{E}_{x\sim q_2}[\ell_t(x)].
\]
La calidad, la diversidad y la tasa de repetición de $q_2$ son variables distintas de la proporción $T_2/(T_1+T_2)$; solo al cambiar una a la vez puede determinarse si la ganancia proviene de la blend o de la duration.


\lecturefigure{slide-049.jpg}{Two-phase approach: phase 1 de amplia cobertura y phase 2 de alta calidad}{CS25 V5 official deck, p.~49}


\lecturefigure{slide-050.jpg}{Data sources: Common Crawl, Wikipedia, libros, artículos y datos de tareas}{CS25 V5 official deck, p.~50}


\lecturefigure{slide-051.jpg}{Evaluation suite: razonamiento, conocimiento, matemáticas, programación y sentido común}{CS25 V5 official deck, p.~51}


\lecturefigure{slide-052.jpg}{Model/training setup: arquitectura fija y configuración de entrenamiento uniforme}{CS25 V5 official deck, p.~52}

\begin{knowledgebox}{Upsampling y effective repetition}
Si una fuente de datos tiene originalmente $N$ tokens y durante el entrenamiento se muestrean $rN$, el upsampling factor es $r$. Un valor $r>1$ aumenta el peso del gradiente de ese dominio, pero también incrementa el riesgo de memorization y de diminishing return. El volumen de datos, la probabilidad de muestreo y el número de pasos de entrenamiento deben reportarse juntos.
\end{knowledgebox}

\subsection{Blend strategy: alta calidad no equivale a una sola fuente}

La formalización definió $q_2$; ahora hay que preguntar de qué sources se compone. Lo que se llama high-quality no es el nombre de una fuente, sino una combinación de filtrado, estructura del documento, densidad de hechos, correspondencia con el dominio y tasa de repetición; un libro de alta calidad y un fragmento de código de alta calidad también contribuyen de manera distinta a cada benchmark. Por ello, al leer una tabla de blends no basta con mirar la quality label: hay que rastrear la proporción de cada categoría, el número real de repeticiones de los tokens de cada categoría y si la mejora aparece a la vez en conocimiento, razonamiento, matemáticas y code.

La segunda fase vuelve a mezclar web high-, medium- y low-quality con curated sources. Al leer una tabla de blends, primero se observa la proporción de cada tipo de datos y luego si las métricas de las tareas mejoran de forma simultánea: una blend puede mejorar el conocimiento pero perjudicar el code, o elevar el promedio pero reducir la cobertura de la cola larga.


\lecturefigure{slide-053.jpg}{Two-phase data blending: phase 1 de amplia cobertura, phase 2 orientada a alta calidad}{CS25 V5 official deck, p.~53}


\lecturefigure{slide-054.jpg}{Impacto de la Phase 2 blending en el rendimiento: resultados multitarea de distintas recetas}{CS25 V5 official deck, p.~54}

\begin{knowledgebox}{Lectura a fondo: cómo evitar la average-score illusion en las tablas de Phase-2}
El promedio de una tabla multitarea oculta los conflictos entre tareas. Primero se agrupan las métricas en knowledge, reasoning, math, code y commonsense, y se observa si una blend produce un desplazamiento general o solo mejora en un grupo; después se compara la absolute gain con la baseline variance, para no presentar una fluctuación mínima como una mejora estable. Si la receta incluye más task-like curated data, también hay que revisar la benchmark format contamination. Las gráficas de pastel representan la sampling share, no la unique data share: una fuente de datos con una proporción pequeña pero muestreada repetidamente puede aportar una gran cantidad de gradiente. Por último, se dividen las gains entre el compute adicional y el data curation cost para obtener la ganancia por unidad de costo. Solo cuando la ventaja se mantiene en varios grupos de tareas, con varias seeds y a mayor escala, hay fundamentos para afirmar que la schedule mejoró el general pretraining, y no que se hizo un ajuste local a la evaluación.
\end{knowledgebox}

\teachervoice{00:22:30--00:23:11, el ponente dice que todas las phase-two blends probadas superan a la simple continuation y mantienen la ventaja al aumentar tokens/modelo. Las notas conservan el alcance de «tested blends» y no lo extrapolan a que cualquier receta de dos fases sea eficaz.}

\subsection{Sobremuestreo y duración de la fase: más no siempre es mejor}

Una vez definida la blend, el entrenamiento todavía enfrenta una cuestión temporal: cuántas veces deben aparecer los mismos tokens de alta calidad y en qué longitud de la fase final deben concentrarse. Si solo se observa la final loss, es fácil confundir la memorization con aprendizaje efectivo; una lectura más confiable considera a la vez la held-out task, la domain diversity y la phase length curve. El turning point de la figura no es una constante fija, sino un recordatorio de que el equipo debe buscar la duration en lugar de suponer que «cuanto más largo, mejor».

Si la segunda fase muestra muchas veces una pequeña cantidad de datos de alta calidad, la loss de entrenamiento puede seguir bajando, pero la ganancia en generalización se satura. Este compromiso puede expresarse de forma simplificada así:
\[
G(r)=B(r)-\lambda M(r),
\]
donde $B(r)$ es la ganancia que aporta la mayor exposición a datos de alta calidad, $M(r)$ es el costo de memorización y estrechamiento de la distribución causado por la repetición, y $\lambda$ representa la sensibilidad de la tarea al sobreajuste. El $r$ óptimo suele estar en un punto interior y no en el infinito.


\lecturefigure{slide-055.jpg}{¿Es eficaz el Phase 2 upsampling?: ganancia, proporción y rendimientos decrecientes}{CS25 V5 official deck, p.~55}


\lecturefigure{slide-056.jpg}{Two-phase scalability: tendencias con mayor token budget y modelos más grandes}{CS25 V5 official deck, p.~56}


\lecturefigure{slide-057.jpg}{Phase duration: una fase final de alta calidad demasiado corta o demasiado larga puede no ser óptima}{CS25 V5 official deck, p.~57}

\begin{knowledgebox}{Lectura a fondo: por qué la phase duration cambia el modelo y no solo el peso de los datos}
Al inicio del entrenamiento, las representaciones del modelo todavía se están formando, y los datos de amplia cobertura permiten construir vocabulario, hechos y patrones generales; en la etapa final, la learning rate, el optimizer state y la basin en la que se encuentran los parámetros ya son distintos, por lo que el mismo lote de tokens de alta calidad produce en ese momento direcciones de actualización diferentes. Por ello, la two-phase schedule no es una simple ponderación de dos conjuntos de datos estáticos: también aprovecha la dependencia de la trayectoria de optimización. Al leer la curva de duration hay que reportar también la learning-rate schedule y si se reinicia el optimizer; si la phase 2 coincide con un cambio en la tasa de aprendizaje, la ganancia puede deberse a ambos factores. También hay que revisar si las tareas que mejora la phase 2 sacrifican capacidades multilingual, creative o long-tail. En la práctica de ingeniería, la phase duration puede plantearse como una Pareto search: en el eje horizontal, los unique tokens/repetitions de alta calidad; en el eje vertical, la ganancia multitarea y el olvido; después se elige una región estable y no un único punto máximo.
\end{knowledgebox}

\teachervoice{00:23:11--00:24:00, en clase se enfatiza que la phase 2 no puede alargarse indefinidamente y que un upsampling excesivo produce detrimental o diminishing returns. Esta caveat se acerca más a una receta aplicable que la afirmación de que «los datos de alta calidad son eficaces».}

\begin{lstlisting}[caption={Esqueleto de experimento reproducible para una two-phase data schedule}]
phase1 = mix(broad_web, books, code, multilingual)
phase2_candidates = propose_quality_blends()
for blend in phase2_candidates:
    proxy = train_small_model(phase1, blend, fixed_tokens=True)
    evaluate(proxy, multi_task_suite, multiple_seeds=True)
selected = choose_pareto_blend(quality, diversity, stability)
train_full_model(phase1, selected, report_repetition=True)
\end{lstlisting}

\subsection{Integrar las dos investigaciones en una data strategy}

TinyDialogues modifica la source structure; two-phase pretraining modifica la schedule. Ambas muestran que la data effectiveness no es una puntuación intrínseca y única de cada muestra, sino una relación de correspondencia entre source, model, objective, etapa de entrenamiento y evaluación.


\lecturefigure{slide-058.jpg}{Two-phase findings: conclusión integral sobre datos estructurados de alta calidad y escalabilidad}{CS25 V5 official deck, p.~58}


\lecturefigure{slide-060.jpg}{Unified insights: data lessons comunes a los datos pequeños y al pretraining a gran escala}{CS25 V5 official deck, p.~60}


\lecturefigure{slide-061.jpg}{Síntesis de la investigación sobre datos: calidad, estructura, programación temporal y experimentos reproducibles son igual de importantes}{CS25 V5 official deck, p.~61}

\begin{knowledgebox}{Términos clave: source, blend, schedule y curriculum}
\textbf{Source} es la categoría de datos original; \textbf{blend} es la proporción de muestreo dentro de una misma fase; \textbf{schedule} describe cómo cambia esa proporción a lo largo de los pasos de entrenamiento; \textbf{curriculum} además puede ordenar los datos de forma intencional según dificultad o capacidad. Two-phase pretraining es una piecewise-constant schedule y no equivale a cualquier curriculum. Al leer las slides 53--61 hay que reportar a la vez el número de tokens por phase, los unique tokens de cada tipo de datos, las effective repetitions y la escala del modelo; de lo contrario, no es posible reproducir la intensidad real de la «fase de alta calidad».
\end{knowledgebox}

\begin{importantbox}{Data-centric decision rule}
Primero hay que preguntar «qué capacidad falta» y después elegir la source que aporte la supervisión correspondiente; seleccionar la blend con experimentos a pequeña escala con varias seeds; al aumentar la escala, verificar que el orden se mantenga; por último, reportar la tasa de repetición, los costos de derechos de autor y privacidad, y los de long-tail. No hay que tomar una tabla de aggregate score como una universal recipe.
\end{importantbox}

\subsection{Resumen de la sección}

Lo esencial de two-phase pretraining no es «dar buenos datos al final», sino convertir la blend y la duration en variables que pueden someterse a ablación. Una fase de alta calidad puede mejorar el rendimiento multitarea, pero el sobremuestreo tiene rendimientos decrecientes; la vía de ingeniería más confiable combina proxy experiment, multi-task evaluation y scale validation.
